\chapter{Expectation-Maximization}

\section{Ziel dieses Kapitels}

Im vorherigen Kapitel haben wir Gaussian Mixture Models (GMMs) betrachtet.
Dort treten latente Variablen (latent variables) auf.
Für jeden Datenpunkt gibt es eine unbekannte Komponentenzugehörigkeit

\[
    z_n.
\]

Wenn diese Zugehörigkeiten bekannt wären, könnten wir die Parameter des Modells relativ einfach schätzen.
Da sie aber nicht beobachtet werden, brauchen wir ein Verfahren, das mit fehlenden oder latenten Variablen umgehen kann.

Dieses Verfahren heißt Expectation-Maximization, kurz EM.

\begin{conceptbox}
Expectation-Maximization (EM) ist ein iterativer Algorithmus zur Maximum-Likelihood-Schätzung bei Modellen mit latenten Variablen.

EM wechselt zwischen zwei Schritten:

\[
    \text{E-Schritt}
    \quad
    \longrightarrow
    \quad
    \text{M-Schritt}.
\]

Im E-Schritt werden Erwartungen über latente Variablen berechnet.
Im M-Schritt werden die Modellparameter aktualisiert.
\end{conceptbox}

Nach diesem Kapitel solltest du folgende Fragen beantworten können:

\begin{itemize}
    \item Was sind latente Variablen (latent variables)?
    \item Warum ist Maximum-Likelihood mit latenten Variablen schwierig?
    \item Was ist der Unterschied zwischen vollständigen und unvollständigen Daten?
    \item Was ist die vollständige Daten-Log-Likelihood (complete-data log-likelihood)?
    \item Was ist die beobachtete Daten-Log-Likelihood (observed-data log-likelihood)?
    \item Was passiert im E-Schritt?
    \item Was passiert im M-Schritt?
    \item Was ist die \(Q\)-Funktion?
    \item Warum verringert EM die Log-Likelihood nicht?
    \item Wie hängt EM mit Jensen-Ungleichung und unteren Schranken zusammen?
    \item Wie wird EM auf Gaussian Mixture Models angewendet?
    \item Wie hängt EM mit \(k\)-means zusammen?
    \item Welche Probleme und Grenzen hat EM?
\end{itemize}

\section{Motivation}

In vielen Modellen gibt es Größen, die nicht direkt beobachtet werden.
Diese Größen nennt man latente Variablen (latent variables).

Beispiele:

\begin{itemize}
    \item Bei einem Gaussian Mixture Model ist unbekannt, aus welcher Komponente ein Datenpunkt stammt.
    \item Bei Clustering ist unbekannt, zu welchem Cluster ein Datenpunkt gehört.
    \item Bei Hidden Markov Models sind die verborgenen Zustände nicht direkt beobachtet.
    \item Bei fehlenden Daten sind manche Merkmale nicht beobachtet.
    \item Bei Topic Models sind die Themenzugehörigkeiten von Wörtern latent.
\end{itemize}

\begin{definitionbox}
Eine latente Variable (latent variable) ist eine Variable, die Teil des Modells ist, aber nicht direkt beobachtet wird.
\end{definitionbox}

Wir unterscheiden:

\[
    X
\]

für beobachtete Daten und

\[
    Z
\]

für latente Variablen.

Die Parameter des Modells bezeichnen wir mit

\[
    \theta.
\]

\section{Beobachtete und vollständige Daten}

Angenommen, wir beobachten nur

\[
    X.
\]

Die latenten Variablen

\[
    Z
\]

sind unbekannt.

Dann nennt man \(X\) die beobachteten Daten (observed data).
Das Paar

\[
    (X,Z)
\]

nennt man vollständige Daten (complete data).

\begin{definitionbox}
Die beobachteten Daten (observed data) sind die tatsächlich gemessenen Daten \(X\).

Die vollständigen Daten (complete data) bestehen aus beobachteten und latenten Variablen:

\[
    (X,Z).
\]
\end{definitionbox}

Bei einem Gaussian Mixture Model sind die beobachteten Daten:

\[
    X
    =
    \{\vec{x}_1,\dots,\vec{x}_N\}.
\]

Die latenten Variablen sind die unbekannten Komponentenzugehörigkeiten:

\[
    Z
    =
    \{z_1,\dots,z_N\}.
\]

\section{Beobachtete Daten-Likelihood}

Unser Ziel ist Maximum-Likelihood-Schätzung (maximum likelihood estimation).
Wir möchten also die Parameter \(\theta\) so wählen, dass die beobachteten Daten möglichst wahrscheinlich sind.

Die beobachtete Daten-Likelihood ist:

\[
    p(X\mid\theta).
\]

Da die latenten Variablen \(Z\) nicht beobachtet werden, müssen wir sie aus der gemeinsamen Verteilung herausmarginalisieren.

Im diskreten Fall:

\[
    p(X\mid\theta)
    =
    \sum_Z
    p(X,Z\mid\theta).
\]

Im kontinuierlichen Fall:

\[
    p(X\mid\theta)
    =
    \int
    p(X,Z\mid\theta)
    \,dZ.
\]

\begin{definitionbox}
Die beobachtete Daten-Likelihood (observed-data likelihood) ist:

\[
    p(X\mid\theta)
    =
    \sum_Z
    p(X,Z\mid\theta).
\]

Dabei wird über die latenten Variablen marginalisiert.
\end{definitionbox}

Die beobachtete Daten-Log-Likelihood ist:

\[
    \ell(\theta)
    =
    \log p(X\mid\theta).
\]

Also im diskreten Fall:

\[
    \ell(\theta)
    =
    \log
    \sum_Z
    p(X,Z\mid\theta).
\]

\begin{conceptbox}
Das schwierige Problem bei latenten Variablen ist der Ausdruck

\[
    \log
    \sum_Z
    p(X,Z\mid\theta).
\]

Der Logarithmus steht außerhalb der Summe über latente Variablen.
Das macht direkte Maximierung oft schwierig.
\end{conceptbox}

\section{Vollständige Daten-Likelihood}

Wenn wir \(Z\) kennen würden, könnten wir stattdessen die vollständige Daten-Likelihood betrachten:

\[
    p(X,Z\mid\theta).
\]

Die vollständige Daten-Log-Likelihood ist:

\[
    \log p(X,Z\mid\theta).
\]

\begin{definitionbox}
Die vollständige Daten-Log-Likelihood (complete-data log-likelihood) ist:

\[
    \log p(X,Z\mid\theta).
\]

Sie wäre einfach zu verwenden, wenn die latenten Variablen \(Z\) beobachtet wären.
\end{definitionbox}

Bei vielen Modellen ist

\[
    \log p(X,Z\mid\theta)
\]

viel einfacher als

\[
    \log p(X\mid\theta).
\]

EM nutzt genau diese Tatsache.

\section{Grundidee von EM}

Da \(Z\) unbekannt ist, können wir nicht direkt mit

\[
    \log p(X,Z\mid\theta)
\]

arbeiten.
Stattdessen berechnen wir im E-Schritt eine Verteilung über \(Z\), gegeben die aktuellen Parameter.

Dann maximieren wir im M-Schritt den Erwartungswert der vollständigen Daten-Log-Likelihood.

\begin{conceptbox}
EM ersetzt die unbekannten latenten Variablen nicht durch einzelne feste Werte.

Stattdessen verwendet EM ihre Posterior-Verteilung unter den aktuellen Parametern.
\end{conceptbox}

Seien die aktuellen Parameter

\[
    \theta^{\mathrm{old}}.
\]

Im E-Schritt berechnen wir:

\[
    p(Z\mid X,\theta^{\mathrm{old}}).
\]

Im M-Schritt maximieren wir:

\[
    Q(\theta,\theta^{\mathrm{old}})
    =
    \mathbb{E}_{Z\mid X,\theta^{\mathrm{old}}}
    [
        \log p(X,Z\mid\theta)
    ].
\]

\section{Die \(Q\)-Funktion}

Die zentrale Größe bei EM ist die \(Q\)-Funktion.

\begin{definitionbox}
Die \(Q\)-Funktion ist der Erwartungswert der vollständigen Daten-Log-Likelihood bezüglich der Posterior-Verteilung der latenten Variablen unter den alten Parametern:

\[
    Q(\theta,\theta^{\mathrm{old}})
    =
    \mathbb{E}_{Z\mid X,\theta^{\mathrm{old}}}
    [
        \log p(X,Z\mid\theta)
    ].
\]
\end{definitionbox}

Im diskreten Fall ist:

\[
    Q(\theta,\theta^{\mathrm{old}})
    =
    \sum_Z
    p(Z\mid X,\theta^{\mathrm{old}})
    \log p(X,Z\mid\theta).
\]

Im kontinuierlichen Fall:

\[
    Q(\theta,\theta^{\mathrm{old}})
    =
    \int
    p(Z\mid X,\theta^{\mathrm{old}})
    \log p(X,Z\mid\theta)
    \,dZ.
\]

\begin{notebox}
In \(Q(\theta,\theta^{\mathrm{old}})\) gibt es zwei Parameterrollen:

\begin{itemize}
    \item \(\theta^{\mathrm{old}}\) bestimmt die Verteilung über die latenten Variablen.
    \item \(\theta\) ist der Parameter, bezüglich dessen maximiert wird.
\end{itemize}
\end{notebox}

\section{Der EM-Algorithmus}

EM besteht aus zwei Schritten.

\begin{definitionbox}
Expectation-Maximization-Algorithmus:

\begin{enumerate}
    \item Initialisiere die Parameter \(\theta^{\mathrm{old}}\).
    \item E-Schritt:
    Berechne
    \[
        p(Z\mid X,\theta^{\mathrm{old}})
    \]
    und damit
    \[
        Q(\theta,\theta^{\mathrm{old}}).
    \]
    \item M-Schritt:
    Wähle neue Parameter
    \[
        \theta^{\mathrm{new}}
        =
        \arg\max_{\theta}
        Q(\theta,\theta^{\mathrm{old}}).
    \]
    \item Setze
    \[
        \theta^{\mathrm{old}}
        \leftarrow
        \theta^{\mathrm{new}}.
    \]
    \item Wiederhole, bis die Log-Likelihood kaum noch steigt.
\end{enumerate}
\end{definitionbox}

\begin{conceptbox}
E-Schritt bedeutet:
Berechne Erwartungen über latente Variablen.

M-Schritt bedeutet:
Maximiere die erwartete vollständige Daten-Log-Likelihood.
\end{conceptbox}

\section{Warum EM die Log-Likelihood nicht verringert}

Eine zentrale Eigenschaft von EM ist:

\[
    \ell(\theta^{\mathrm{new}})
    \geq
    \ell(\theta^{\mathrm{old}}).
\]

Die beobachtete Daten-Log-Likelihood sinkt also nicht.

Um das zu verstehen, betrachten wir eine untere Schranke (lower bound) der Log-Likelihood.

\section{Untere Schranke mit einer Hilfsverteilung}

Sei \(q(Z)\) eine beliebige Verteilung über die latenten Variablen.
Dann gilt:

\[
    \log p(X\mid\theta)
    =
    \log
    \sum_Z
    p(X,Z\mid\theta).
\]

Wir multiplizieren und dividieren mit \(q(Z)\):

\[
    \log p(X\mid\theta)
    =
    \log
    \sum_Z
    q(Z)
    \frac{
        p(X,Z\mid\theta)
    }{
        q(Z)
    }.
\]

Da \(q(Z)\) eine Verteilung ist, ist die Summe ein Erwartungswert bezüglich \(q\):

\[
    \log p(X\mid\theta)
    =
    \log
    \mathbb{E}_{q}
    \left[
        \frac{
            p(X,Z\mid\theta)
        }{
            q(Z)
        }
    \right].
\]

Da \(\log\) konkav ist, folgt mit Jensen-Ungleichung:

\[
    \log
    \mathbb{E}_{q}
    \left[
        \frac{
            p(X,Z\mid\theta)
        }{
            q(Z)
        }
    \right]
    \geq
    \mathbb{E}_{q}
    \left[
        \log
        \frac{
            p(X,Z\mid\theta)
        }{
            q(Z)
        }
    \right].
\]

Also:

\[
    \log p(X\mid\theta)
    \geq
    \sum_Z
    q(Z)
    \log
    \frac{
        p(X,Z\mid\theta)
    }{
        q(Z)
    }.
\]

\begin{definitionbox}
Die EM-Untergrenze oder Free Energy ist:

\[
    \mathcal{F}(q,\theta)
    =
    \sum_Z
    q(Z)
    \log
    \frac{
        p(X,Z\mid\theta)
    }{
        q(Z)
    }.
\]

Sie erfüllt:

\[
    \mathcal{F}(q,\theta)
    \leq
    \log p(X\mid\theta).
\]
\end{definitionbox}

\section{Zerlegung der unteren Schranke}

Wir können die Free Energy schreiben als:

\[
    \mathcal{F}(q,\theta)
    =
    \sum_Z
    q(Z)
    \log p(X,Z\mid\theta)
    -
    \sum_Z
    q(Z)
    \log q(Z).
\]

Der erste Term ist der erwartete vollständige Daten-Log-Likelihood-Term.
Der zweite Term ist die Entropie von \(q\):

\[
    H(q)
    =
    -\sum_Z
    q(Z)\log q(Z).
\]

Damit:

\[
    \mathcal{F}(q,\theta)
    =
    \mathbb{E}_{q}
    [
        \log p(X,Z\mid\theta)
    ]
    +
    H(q).
\]

\begin{conceptbox}
Die Free Energy besteht aus:

\[
    \text{erwartete vollständige Daten-Log-Likelihood}
    +
    \text{Entropie der Hilfsverteilung}.
\]
\end{conceptbox}

\section{Zusammenhang mit KL-Divergenz}

Die Differenz zwischen Log-Likelihood und Free Energy ist eine KL-Divergenz (Kullback-Leibler divergence).

Es gilt:

\[
    \log p(X\mid\theta)
    =
    \mathcal{F}(q,\theta)
    +
    \mathrm{KL}
    \left(
        q(Z)
        \,\|\, 
        p(Z\mid X,\theta)
    \right).
\]

Da die KL-Divergenz immer nichtnegativ ist,

\[
    \mathrm{KL}(q\|p)\geq 0,
\]

folgt:

\[
    \mathcal{F}(q,\theta)
    \leq
    \log p(X\mid\theta).
\]

\begin{definitionbox}
Die KL-Divergenz zwischen zwei Verteilungen \(q\) und \(p\) ist:

\[
    \mathrm{KL}(q\|p)
    =
    \sum_Z
    q(Z)
    \log
    \frac{q(Z)}{p(Z)}.
\]

Sie ist immer nichtnegativ.
\end{definitionbox}

Die Schranke ist genau dann eng, wenn

\[
    q(Z)
    =
    p(Z\mid X,\theta).
\]

Dann ist:

\[
    \mathrm{KL}
    \left(
        q(Z)
        \,\|\,
        p(Z\mid X,\theta)
    \right)
    =
    0.
\]

\section{E-Schritt als Optimierung von \(q\)}

Im E-Schritt halten wir \(\theta\) fest bei

\[
    \theta^{\mathrm{old}}.
\]

Wir wählen \(q(Z)\) so, dass die Free Energy maximal wird.
Das passiert, wenn die KL-Divergenz verschwindet.

Also:

\[
    q^{\mathrm{new}}(Z)
    =
    p(Z\mid X,\theta^{\mathrm{old}}).
\]

\begin{definitionbox}
Im E-Schritt wird die Hilfsverteilung auf den Posterior der latenten Variablen gesetzt:

\[
    q(Z)
    =
    p(Z\mid X,\theta^{\mathrm{old}}).
\]
\end{definitionbox}

Damit wird die untere Schranke bei den alten Parametern eng:

\[
    \mathcal{F}(q,\theta^{\mathrm{old}})
    =
    \log p(X\mid\theta^{\mathrm{old}}).
\]

\section{M-Schritt als Optimierung von \(\theta\)}

Im M-Schritt halten wir \(q(Z)\) fest und maximieren die Free Energy bezüglich \(\theta\).

Da der Entropieterm

\[
    H(q)
\]

nicht von \(\theta\) abhängt, maximieren wir nur:

\[
    \mathbb{E}_{q}
    [
        \log p(X,Z\mid\theta)
    ].
\]

Mit

\[
    q(Z)
    =
    p(Z\mid X,\theta^{\mathrm{old}})
\]

ist das genau:

\[
    Q(\theta,\theta^{\mathrm{old}})
    =
    \mathbb{E}_{Z\mid X,\theta^{\mathrm{old}}}
    [
        \log p(X,Z\mid\theta)
    ].
\]

\begin{definitionbox}
Im M-Schritt wird

\[
    Q(\theta,\theta^{\mathrm{old}})
\]

bezüglich \(\theta\) maximiert:

\[
    \theta^{\mathrm{new}}
    =
    \arg\max_{\theta}
    Q(\theta,\theta^{\mathrm{old}}).
\]
\end{definitionbox}

\section{Monotonie von EM}

Nun können wir die Monotonie von EM verstehen.

Nach dem E-Schritt ist die Schranke bei \(\theta^{\mathrm{old}}\) eng:

\[
    \mathcal{F}(q,\theta^{\mathrm{old}})
    =
    \ell(\theta^{\mathrm{old}}).
\]

Im M-Schritt wählen wir \(\theta^{\mathrm{new}}\), sodass:

\[
    \mathcal{F}(q,\theta^{\mathrm{new}})
    \geq
    \mathcal{F}(q,\theta^{\mathrm{old}}).
\]

Da die Free Energy eine untere Schranke der Log-Likelihood ist:

\[
    \ell(\theta^{\mathrm{new}})
    \geq
    \mathcal{F}(q,\theta^{\mathrm{new}}).
\]

Damit:

\[
    \ell(\theta^{\mathrm{new}})
    \geq
    \mathcal{F}(q,\theta^{\mathrm{new}})
    \geq
    \mathcal{F}(q,\theta^{\mathrm{old}})
    =
    \ell(\theta^{\mathrm{old}}).
\]

Also:

\[
    \ell(\theta^{\mathrm{new}})
    \geq
    \ell(\theta^{\mathrm{old}}).
\]

\begin{conceptbox}
EM verringert die beobachtete Daten-Log-Likelihood nicht.

Das bedeutet aber nicht, dass EM das globale Maximum findet.
EM kann in lokalen Maxima oder Sattelpunkten landen.
\end{conceptbox}

\section{EM für Gaussian Mixture Models}

Nun wenden wir EM auf Gaussian Mixture Models an.

Das Modell ist:

\[
    p(\vec{x})
    =
    \sum_{k=1}^{K}
    \pi_k
    \mathcal{N}(\vec{x}\mid\vec{\mu}_k,\Sigma_k).
\]

Die Parameter sind:

\[
    \Theta
    =
    \{\pi_k,\vec{\mu}_k,\Sigma_k\}_{k=1}^{K}.
\]

Die latente Variable \(z_n\) gibt an, aus welcher Komponente Datenpunkt \(\vec{x}_n\) stammt.

\section{Vollständige Daten-Likelihood für GMMs}

Wir verwenden One-Hot-Variablen:

\[
    z_{nk}\in\{0,1\},
\]

mit

\[
    \sum_{k=1}^{K}
    z_{nk}=1.
\]

Die Verteilung der latenten Variablen ist:

\[
    p(\vec{z}_n\mid\pi)
    =
    \prod_{k=1}^{K}
    \pi_k^{z_{nk}}.
\]

Die bedingte Verteilung des Datenpunkts ist:

\[
    p(\vec{x}_n\mid\vec{z}_n,\Theta)
    =
    \prod_{k=1}^{K}
    \mathcal{N}(\vec{x}_n\mid\vec{\mu}_k,\Sigma_k)^{z_{nk}}.
\]

Die gemeinsame Verteilung ist:

\[
    p(\vec{x}_n,\vec{z}_n\mid\Theta)
    =
    \prod_{k=1}^{K}
    \left[
        \pi_k
        \mathcal{N}(\vec{x}_n\mid\vec{\mu}_k,\Sigma_k)
    \right]^{z_{nk}}.
\]

Für alle Datenpunkte:

\[
    p(X,Z\mid\Theta)
    =
    \prod_{n=1}^{N}
    \prod_{k=1}^{K}
    \left[
        \pi_k
        \mathcal{N}(\vec{x}_n\mid\vec{\mu}_k,\Sigma_k)
    \right]^{z_{nk}}.
\]

\begin{definitionbox}
Die vollständige Daten-Likelihood eines GMMs ist:

\[
    p(X,Z\mid\Theta)
    =
    \prod_{n=1}^{N}
    \prod_{k=1}^{K}
    \left[
        \pi_k
        \mathcal{N}(\vec{x}_n\mid\vec{\mu}_k,\Sigma_k)
    \right]^{z_{nk}}.
\]
\end{definitionbox}

\section{Vollständige Daten-Log-Likelihood für GMMs}

Wir nehmen den Logarithmus:

\[
    \log p(X,Z\mid\Theta)
    =
    \log
    \prod_{n=1}^{N}
    \prod_{k=1}^{K}
    \left[
        \pi_k
        \mathcal{N}(\vec{x}_n\mid\vec{\mu}_k,\Sigma_k)
    \right]^{z_{nk}}.
\]

Der Logarithmus macht aus Produkten Summen:

\[
    \log p(X,Z\mid\Theta)
    =
    \sum_{n=1}^{N}
    \sum_{k=1}^{K}
    z_{nk}
    \log
    \left[
        \pi_k
        \mathcal{N}(\vec{x}_n\mid\vec{\mu}_k,\Sigma_k)
    \right].
\]

Mit

\[
    \log(ab)=\log a+\log b
\]

ergibt sich:

\[
    \log p(X,Z\mid\Theta)
    =
    \sum_{n=1}^{N}
    \sum_{k=1}^{K}
    z_{nk}
    \left[
        \log \pi_k
        +
        \log
        \mathcal{N}(\vec{x}_n\mid\vec{\mu}_k,\Sigma_k)
    \right].
\]

\begin{definitionbox}
Die vollständige Daten-Log-Likelihood eines GMMs ist:

\[
    \log p(X,Z\mid\Theta)
    =
    \sum_{n=1}^{N}
    \sum_{k=1}^{K}
    z_{nk}
    \left[
        \log \pi_k
        +
        \log
        \mathcal{N}(\vec{x}_n\mid\vec{\mu}_k,\Sigma_k)
    \right].
\]
\end{definitionbox}

\section{E-Schritt für GMMs}

Im E-Schritt berechnen wir den Erwartungswert von \(z_{nk}\) unter dem Posterior:

\[
    \mathbb{E}[z_{nk}]
    =
    p(z_{nk}=1\mid\vec{x}_n,\Theta^{\mathrm{old}}).
\]

Diese Größe ist die Responsibility:

\[
    \gamma_{nk}
    =
    p(z_n=k\mid\vec{x}_n,\Theta^{\mathrm{old}}).
\]

Mit Bayes' Regel:

\[
    \gamma_{nk}
    =
    \frac{
        \pi_k^{\mathrm{old}}
        \mathcal{N}(\vec{x}_n\mid\vec{\mu}_k^{\mathrm{old}},\Sigma_k^{\mathrm{old}})
    }{
        \sum_{j=1}^{K}
        \pi_j^{\mathrm{old}}
        \mathcal{N}(\vec{x}_n\mid\vec{\mu}_j^{\mathrm{old}},\Sigma_j^{\mathrm{old}})
    }.
\]

\begin{definitionbox}
Der E-Schritt für ein GMM berechnet:

\[
    \gamma_{nk}
    =
    \frac{
        \pi_k
        \mathcal{N}(\vec{x}_n\mid\vec{\mu}_k,\Sigma_k)
    }{
        \sum_{j=1}^{K}
        \pi_j
        \mathcal{N}(\vec{x}_n\mid\vec{\mu}_j,\Sigma_j)
    }.
\]

Dabei werden die aktuellen Parameter verwendet.
\end{definitionbox}

Für jeden Datenpunkt gilt:

\[
    \sum_{k=1}^{K}
    \gamma_{nk}=1.
\]

\section{\(Q\)-Funktion für GMMs}

Die \(Q\)-Funktion ist der Erwartungswert der vollständigen Daten-Log-Likelihood.

Da

\[
    \mathbb{E}[z_{nk}]
    =
    \gamma_{nk},
\]

erhalten wir:

\[
    Q(\Theta,\Theta^{\mathrm{old}})
    =
    \sum_{n=1}^{N}
    \sum_{k=1}^{K}
    \gamma_{nk}
    \left[
        \log \pi_k
        +
        \log
        \mathcal{N}(\vec{x}_n\mid\vec{\mu}_k,\Sigma_k)
    \right].
\]

\begin{definitionbox}
Die \(Q\)-Funktion für ein GMM ist:

\[
    Q(\Theta,\Theta^{\mathrm{old}})
    =
    \sum_{n=1}^{N}
    \sum_{k=1}^{K}
    \gamma_{nk}
    \left[
        \log \pi_k
        +
        \log
        \mathcal{N}(\vec{x}_n\mid\vec{\mu}_k,\Sigma_k)
    \right].
\]
\end{definitionbox}

Im M-Schritt maximieren wir diese Funktion bezüglich

\[
    \pi_k,
    \qquad
    \vec{\mu}_k,
    \qquad
    \Sigma_k.
\]

\section{M-Schritt: effektive Clustergröße}

Zuerst definieren wir die effektive Clustergröße:

\[
    N_k
    =
    \sum_{n=1}^{N}
    \gamma_{nk}.
\]

\begin{definitionbox}
Die effektive Anzahl der Punkte in Komponente \(k\) ist:

\[
    N_k
    =
    \sum_{n=1}^{N}
    \gamma_{nk}.
\]
\end{definitionbox}

Bei harten Zuordnungen wäre \(N_k\) einfach die Anzahl der Punkte im Cluster.
Bei weichen Zuordnungen zählt jeder Punkt anteilig.

\section{M-Schritt: Update für Mittelwerte}

Wir maximieren \(Q\) bezüglich \(\vec{\mu}_k\).
Nur der Gauß-Term hängt von \(\vec{\mu}_k\) ab.

Für eine multivariate Gauß-Verteilung gilt:

\[
    \log
    \mathcal{N}(\vec{x}_n\mid\vec{\mu}_k,\Sigma_k)
    =
    -\frac{D}{2}\log(2\pi)
    -
    \frac{1}{2}\log|\Sigma_k|
    -
    \frac{1}{2}
    (\vec{x}_n-\vec{\mu}_k)^{\top}
    \Sigma_k^{-1}
    (\vec{x}_n-\vec{\mu}_k).
\]

Die Terme, die von \(\vec{\mu}_k\) abhängen, sind:

\[
    -\frac{1}{2}
    \sum_{n=1}^{N}
    \gamma_{nk}
    (\vec{x}_n-\vec{\mu}_k)^{\top}
    \Sigma_k^{-1}
    (\vec{x}_n-\vec{\mu}_k).
\]

Maximieren entspricht Minimieren von:

\[
    \sum_{n=1}^{N}
    \gamma_{nk}
    (\vec{x}_n-\vec{\mu}_k)^{\top}
    \Sigma_k^{-1}
    (\vec{x}_n-\vec{\mu}_k).
\]

Die Lösung ist der gewichtete Mittelwert:

\[
    \vec{\mu}_k^{\mathrm{new}}
    =
    \frac{1}{N_k}
    \sum_{n=1}^{N}
    \gamma_{nk}
    \vec{x}_n.
\]

\begin{definitionbox}
Das EM-Update für den Mittelwert einer GMM-Komponente ist:

\[
    \vec{\mu}_k^{\mathrm{new}}
    =
    \frac{1}{N_k}
    \sum_{n=1}^{N}
    \gamma_{nk}
    \vec{x}_n.
\]
\end{definitionbox}

\section{M-Schritt: Update für Kovarianzen}

Das Update für die Kovarianzmatrix ist die gewichtete empirische Kovarianz:

\[
    \Sigma_k^{\mathrm{new}}
    =
    \frac{1}{N_k}
    \sum_{n=1}^{N}
    \gamma_{nk}
    (\vec{x}_n-\vec{\mu}_k^{\mathrm{new}})
    (\vec{x}_n-\vec{\mu}_k^{\mathrm{new}})^{\top}.
\]

\begin{definitionbox}
Das EM-Update für die Kovarianzmatrix einer GMM-Komponente ist:

\[
    \Sigma_k^{\mathrm{new}}
    =
    \frac{1}{N_k}
    \sum_{n=1}^{N}
    \gamma_{nk}
    (\vec{x}_n-\vec{\mu}_k^{\mathrm{new}})
    (\vec{x}_n-\vec{\mu}_k^{\mathrm{new}})^{\top}.
\]
\end{definitionbox}

\begin{conceptbox}
Die Kovarianz jeder Komponente wird aus allen Datenpunkten berechnet, aber jeder Datenpunkt wird mit seiner Responsibility \(\gamma_{nk}\) gewichtet.
\end{conceptbox}

\section{M-Schritt: Update für Mischgewichte}

Die Mischgewichte müssen die Nebenbedingungen erfüllen:

\[
    \pi_k\geq 0,
    \qquad
    \sum_{k=1}^{K}
    \pi_k=1.
\]

Der relevante Teil der \(Q\)-Funktion ist:

\[
    \sum_{n=1}^{N}
    \sum_{k=1}^{K}
    \gamma_{nk}
    \log \pi_k.
\]

Wir verwenden einen Lagrange-Multiplikator \(\lambda\) für die Nebenbedingung:

\[
    \sum_{k=1}^{K}\pi_k=1.
\]

Die Lagrange-Funktion ist:

\[
    \mathcal{L}
    =
    \sum_{n=1}^{N}
    \sum_{k=1}^{K}
    \gamma_{nk}
    \log \pi_k
    +
    \lambda
    \left(
        \sum_{k=1}^{K}
        \pi_k
        -
        1
    \right).
\]

Ableiten nach \(\pi_k\):

\[
    \frac{\partial\mathcal{L}}{\partial\pi_k}
    =
    \frac{
        \sum_{n=1}^{N}
        \gamma_{nk}
    }{
        \pi_k
    }
    +
    \lambda.
\]

Setze gleich \(0\):

\[
    \frac{N_k}{\pi_k}
    +
    \lambda
    =
    0.
\]

Damit:

\[
    \pi_k
    =
    -\frac{N_k}{\lambda}.
\]

Nun summieren wir über \(k\):

\[
    \sum_{k=1}^{K}
    \pi_k
    =
    -\frac{1}{\lambda}
    \sum_{k=1}^{K}
    N_k.
\]

Da

\[
    \sum_{k=1}^{K}\pi_k=1
\]

und

\[
    \sum_{k=1}^{K}N_k
    =
    N,
\]

folgt:

\[
    1
    =
    -\frac{N}{\lambda}.
\]

Also:

\[
    \lambda=-N.
\]

Damit:

\[
    \pi_k^{\mathrm{new}}
    =
    \frac{N_k}{N}.
\]

\begin{definitionbox}
Das EM-Update für das Mischgewicht ist:

\[
    \pi_k^{\mathrm{new}}
    =
    \frac{N_k}{N}.
\]
\end{definitionbox}

\section{Zusammenfassung der GMM-EM-Updates}

\begin{conceptbox}
Für ein Gaussian Mixture Model lauten die EM-Schritte:

E-Schritt:

\[
    \gamma_{nk}
    =
    \frac{
        \pi_k
        \mathcal{N}(\vec{x}_n\mid\vec{\mu}_k,\Sigma_k)
    }{
        \sum_{j=1}^{K}
        \pi_j
        \mathcal{N}(\vec{x}_n\mid\vec{\mu}_j,\Sigma_j)
    }.
\]

M-Schritt:

\[
    N_k
    =
    \sum_{n=1}^{N}
    \gamma_{nk},
\]

\[
    \vec{\mu}_k
    =
    \frac{1}{N_k}
    \sum_{n=1}^{N}
    \gamma_{nk}
    \vec{x}_n,
\]

\[
    \Sigma_k
    =
    \frac{1}{N_k}
    \sum_{n=1}^{N}
    \gamma_{nk}
    (\vec{x}_n-\vec{\mu}_k)
    (\vec{x}_n-\vec{\mu}_k)^{\top},
\]

\[
    \pi_k
    =
    \frac{N_k}{N}.
\]
\end{conceptbox}

\section{EM und \(k\)-means}

EM für GMMs ist eng verwandt mit \(k\)-means.

Bei \(k\)-means gibt es harte Zuordnungen:

\[
    r_{nk}
    \in
    \{0,1\}.
\]

Bei GMM-EM gibt es weiche Responsibilities:

\[
    \gamma_{nk}
    \in
    [0,1].
\]

Wenn ein GMM gleiche isotrope Kovarianzen hat,

\[
    \Sigma_k=\sigma^2 I,
\]

und

\[
    \sigma^2
    \to
    0,
\]

werden die Responsibilities fast hart.
Dann verhält sich EM ähnlich wie \(k\)-means.

\begin{conceptbox}
\(k\)-means kann als Grenzfall von EM für GMMs verstanden werden:

\[
    \text{weiche Zuordnung}
    \quad
    \longrightarrow
    \quad
    \text{harte Zuordnung}.
\]
\end{conceptbox}

\section{Generalized EM}

Beim klassischen EM wird im M-Schritt \(Q\) vollständig maximiert:

\[
    \theta^{\mathrm{new}}
    =
    \arg\max_{\theta}
    Q(\theta,\theta^{\mathrm{old}}).
\]

Manchmal ist das vollständige Maximum schwer zu finden.
Dann reicht es, \(Q\) nur zu verbessern:

\[
    Q(\theta^{\mathrm{new}},\theta^{\mathrm{old}})
    \geq
    Q(\theta^{\mathrm{old}},\theta^{\mathrm{old}}).
\]

Das nennt man Generalized EM.

\begin{definitionbox}
Generalized EM (GEM) verlangt im M-Schritt nicht das exakte Maximum von \(Q\), sondern nur eine Verbesserung von \(Q\).
\end{definitionbox}

\begin{notebox}
Auch Generalized EM kann die Log-Likelihood monoton verbessern, solange der M-Schritt die untere Schranke erhöht.
\end{notebox}

\section{Hard EM}

Eine weitere Variante ist Hard EM.
Dabei verwendet man nicht die vollständige Posterior-Verteilung der latenten Variablen.
Stattdessen wählt man die wahrscheinlichste latente Zuordnung.

\[
    z_n^{*}
    =
    \arg\max_k
    p(z_n=k\mid\vec{x}_n,\theta).
\]

Dann wird so getan, als wären diese Zuordnungen bekannt.

\begin{definitionbox}
Hard EM ersetzt weiche Posterior-Zuordnungen durch harte Zuordnungen:

\[
    z_n^{*}
    =
    \arg\max_k
    p(z_n=k\mid\vec{x}_n,\theta).
\]
\end{definitionbox}

\begin{conceptbox}
Hard EM ist einfacher, verliert aber Unsicherheitsinformation.

\(k\)-means kann als eine Form von Hard EM für ein stark eingeschränktes GMM verstanden werden.
\end{conceptbox}

\section{EM bei fehlenden Daten}

EM kann auch für fehlende Daten verwendet werden.
Angenommen, ein Datenpunkt besteht aus beobachteten und fehlenden Teilen:

\[
    \vec{x}_n
    =
    (\vec{x}_{n,\mathrm{obs}},\vec{x}_{n,\mathrm{miss}}).
\]

Dann können die fehlenden Werte als latente Variablen behandelt werden.

Im E-Schritt berechnet man Erwartungen über die fehlenden Werte gegeben die beobachteten Werte und aktuelle Parameter:

\[
    p(\vec{x}_{n,\mathrm{miss}}\mid\vec{x}_{n,\mathrm{obs}},\theta^{\mathrm{old}}).
\]

Im M-Schritt aktualisiert man die Parameter basierend auf diesen Erwartungen.

\begin{conceptbox}
EM ist nicht nur für Clustering nützlich.

EM ist ein allgemeines Prinzip für Maximum-Likelihood-Schätzung mit nicht beobachteten oder fehlenden Größen.
\end{conceptbox}

\section{Initialisierung}

EM benötigt Startwerte für die Parameter.
Da EM nur zu lokalen Maxima konvergiert, ist Initialisierung wichtig.

Typische Strategien:

\begin{itemize}
    \item zufällige Initialisierung,
    \item \(k\)-means zur Initialisierung von GMMs,
    \item mehrere Starts und Auswahl der besten Log-Likelihood,
    \item fachlich motivierte Startwerte.
\end{itemize}

\begin{examtipbox}
EM findet im Allgemeinen kein globales Maximum.

Deshalb sollte man bei nichtkonvexen Modellen mehrere Initialisierungen ausprobieren.
\end{examtipbox}

\section{Konvergenzkriterien}

EM wird iteriert, bis ein Stoppkriterium erfüllt ist.

Typische Kriterien:

\begin{itemize}
    \item Änderung der Log-Likelihood ist klein:
    \[
        |\ell^{(\tau+1)}-\ell^{(\tau)}|<\varepsilon.
    \]

    \item Relative Änderung der Log-Likelihood ist klein:
    \[
        \frac{
            |\ell^{(\tau+1)}-\ell^{(\tau)}|
        }{
            |\ell^{(\tau)}|
        }
        <
        \varepsilon.
    \]

    \item Änderung der Parameter ist klein:
    \[
        \|\theta^{(\tau+1)}-\theta^{(\tau)}\|<\varepsilon.
    \]

    \item Maximale Anzahl an Iterationen erreicht.
\end{itemize}

\begin{notebox}
In der Praxis verwendet man oft eine Kombination aus Log-Likelihood-Toleranz und maximaler Iterationszahl.
\end{notebox}

\section{Numerische Stabilität}

Bei EM für GMMs können sehr kleine Wahrscheinlichkeiten auftreten.
Deshalb arbeitet man oft im Log-Raum.

Für die Log-Likelihood:

\[
    \ell(\Theta)
    =
    \sum_{n=1}^{N}
    \log
    \left(
        \sum_{k=1}^{K}
        \pi_k
        \mathcal{N}(\vec{x}_n\mid\vec{\mu}_k,\Sigma_k)
    \right)
\]

verwendet man die Log-Sum-Exp-Technik.

\[
    \log\sum_k\exp(a_k)
    =
    m
    +
    \log
    \sum_k
    \exp(a_k-m),
\]

wobei

\[
    m=\max_k a_k.
\]

\begin{definitionbox}
Log-Sum-Exp stabilisiert Berechnungen der Form

\[
    \log\sum_k\exp(a_k).
\]

Durch Subtraktion des Maximums werden numerische Überläufe vermieden.
\end{definitionbox}

\section{Singularitäten bei GMMs}

Bei GMMs kann EM in degenerierte Lösungen laufen.
Eine Komponente kann sich auf einen einzelnen Datenpunkt konzentrieren und ihre Kovarianz gegen \(0\) schrumpfen lassen.

Dann kann die Likelihood gegen unendlich gehen.

\[
    \Sigma_k\to 0
    \quad
    \Rightarrow
    \quad
    \mathcal{N}(\vec{x}_n\mid\vec{\mu}_k,\Sigma_k)\to\infty
\]

für

\[
    \vec{\mu}_k=\vec{x}_n.
\]

\begin{examtipbox}
Die GMM-Likelihood ist nicht immer wohldefiniert maximierbar, weil singuläre Lösungen auftreten können.
\end{examtipbox}

Gegenmaßnahmen:

\begin{itemize}
    \item Kovarianzregularisierung:
    \[
        \Sigma_k
        \leftarrow
        \Sigma_k+\varepsilon I.
    \]

    \item minimale Varianz erzwingen,
    \item zu kleine Komponenten entfernen,
    \item Anzahl der Komponenten \(K\) reduzieren,
    \item diagonale oder geteilte Kovarianzmatrizen verwenden.
\end{itemize}

\section{Vorteile von EM}

\begin{itemize}
    \item EM nutzt latente Variablen systematisch.
    \item Viele M-Schritte haben geschlossene Formeln.
    \item Die Log-Likelihood wird monoton verbessert.
    \item EM ist intuitiv: fehlende Information schätzen, dann Parameter aktualisieren.
    \item EM ist auf viele Modelle anwendbar.
\end{itemize}

\begin{conceptbox}
EM ist besonders nützlich, wenn die vollständige Daten-Likelihood einfach wäre, aber ein Teil der Daten latent oder fehlend ist.
\end{conceptbox}

\section{Nachteile von EM}

\begin{itemize}
    \item EM findet nur lokale Optima.
    \item EM kann langsam konvergieren.
    \item Initialisierung ist wichtig.
    \item Der E-Schritt kann in komplexen Modellen schwierig sein.
    \item Der M-Schritt kann in komplexen Modellen schwierig sein.
    \item Bei GMMs können singuläre Lösungen auftreten.
\end{itemize}

\begin{notebox}
EM ist kein einzelner spezieller Algorithmus nur für GMMs.
EM ist ein allgemeines Optimierungsprinzip für Modelle mit latenten Variablen.
\end{notebox}

\section{Mini-Aufgaben}

\subsection{Aufgabe 1: Latente Variable erkennen}

In einem Clustering-Modell beobachtet man Datenpunkte \(\vec{x}_n\), aber nicht, zu welchem Cluster sie gehören.
Was ist die latente Variable?

\begin{examplebox}
Die latente Variable ist die Clusterzugehörigkeit:

\[
    z_n.
\]

Sie gibt an, aus welchem Cluster oder welcher Komponente der Datenpunkt stammt.
\end{examplebox}

\subsection{Aufgabe 2: Responsibility berechnen}

Ein GMM hat zwei Komponenten.
Für einen Datenpunkt gelten:

\[
    \pi_1\mathcal{N}_1(\vec{x})=0{,}2,
\]

und

\[
    \pi_2\mathcal{N}_2(\vec{x})=0{,}8.
\]

Berechne \(\gamma_{1}\) und \(\gamma_{2}\).

\begin{examplebox}
Die Summe ist:

\[
    0{,}2+0{,}8=1.
\]

Also:

\[
    \gamma_1
    =
    \frac{0{,}2}{1}
    =
    0{,}2,
\]

und

\[
    \gamma_2
    =
    \frac{0{,}8}{1}
    =
    0{,}8.
\]
\end{examplebox}

\subsection{Aufgabe 3: Effektive Clustergröße}

Für eine Komponente seien die Responsibilities:

\[
    0{,}1,
    \quad
    0{,}4,
    \quad
    0{,}8,
    \quad
    0{,}7.
\]

Berechne \(N_k\).

\begin{examplebox}
Die effektive Clustergröße ist:

\[
    N_k
    =
    \sum_{n=1}^{N}
    \gamma_{nk}.
\]

Also:

\[
    N_k
    =
    0{,}1+0{,}4+0{,}8+0{,}7
    =
    2{,}0.
\]
\end{examplebox}

\subsection{Aufgabe 4: Mischgewicht}

Ein Datensatz hat

\[
    N=20
\]

Datenpunkte.
Für Komponente \(k\) gilt:

\[
    N_k=5.
\]

Berechne das neue Mischgewicht.

\begin{examplebox}
Im M-Schritt gilt:

\[
    \pi_k
    =
    \frac{N_k}{N}.
\]

Einsetzen:

\[
    \pi_k
    =
    \frac{5}{20}
    =
    0{,}25.
\]
\end{examplebox}

\subsection{Aufgabe 5: Monotonie}

Warum sinkt die Log-Likelihood bei EM nicht?

\begin{examplebox}
EM konstruiert eine untere Schranke an die Log-Likelihood.

Im E-Schritt wird diese Schranke bei den aktuellen Parametern eng gemacht.
Im M-Schritt wird die Schranke maximiert oder zumindest erhöht.

Dadurch gilt:

\[
    \ell(\theta^{\mathrm{new}})
    \geq
    \ell(\theta^{\mathrm{old}}).
\]
\end{examplebox}

\section{Häufige Fehler}

\begin{examtipbox}
Typische Fehler bei Expectation-Maximization:

\begin{enumerate}
    \item EM nur als GMM-Algorithmus ansehen.
    \item Latente Variablen mit Parametern verwechseln.
    \item Beobachtete Daten-Likelihood und vollständige Daten-Likelihood verwechseln.
    \item Den Logarithmus falsch durch die Summe ziehen.
    \item Im E-Schritt harte Zuordnungen verwenden, obwohl Soft EM gefragt ist.
    \item Die \(Q\)-Funktion mit der beobachteten Log-Likelihood verwechseln.
    \item Im M-Schritt die Responsibilities nicht als Gewichte verwenden.
    \item \(N_k=\sum_n\gamma_{nk}\) vergessen.
    \item Denken, EM finde garantiert das globale Maximum.
    \item Ignorieren, dass Initialisierung wichtig ist.
    \item Singuläre Kovarianzen bei GMMs nicht beachten.
    \item Die Monotonie von EM mit globaler Optimalität verwechseln.
\end{enumerate}
\end{examtipbox}

\section{Zusammenfassung}

\begin{itemize}
    \item EM ist ein Algorithmus für Maximum-Likelihood-Schätzung mit latenten Variablen.
    \item Beobachtete Daten sind \(X\).
    \item Latente Variablen sind \(Z\).
    \item Die beobachtete Daten-Likelihood ist:
    \[
        p(X\mid\theta)
        =
        \sum_Z
        p(X,Z\mid\theta).
    \]
    \item Die beobachtete Log-Likelihood enthält oft:
    \[
        \log
        \sum_Z
        p(X,Z\mid\theta).
    \]
    \item Die vollständige Daten-Log-Likelihood ist:
    \[
        \log p(X,Z\mid\theta).
    \]
    \item Die \(Q\)-Funktion ist:
    \[
        Q(\theta,\theta^{\mathrm{old}})
        =
        \mathbb{E}_{Z\mid X,\theta^{\mathrm{old}}}
        [
            \log p(X,Z\mid\theta)
        ].
    \]
    \item Im E-Schritt berechnet man:
    \[
        p(Z\mid X,\theta^{\mathrm{old}}).
    \]
    \item Im M-Schritt berechnet man:
    \[
        \theta^{\mathrm{new}}
        =
        \arg\max_{\theta}
        Q(\theta,\theta^{\mathrm{old}}).
    \]
    \item EM kann über eine untere Schranke an die Log-Likelihood hergeleitet werden.
    \item Diese Schranke hängt mit Jensen-Ungleichung und KL-Divergenz zusammen.
    \item EM verringert die Log-Likelihood nicht:
    \[
        \ell(\theta^{\mathrm{new}})
        \geq
        \ell(\theta^{\mathrm{old}}).
    \]
    \item Für GMMs berechnet der E-Schritt Responsibilities:
    \[
        \gamma_{nk}
        =
        \frac{
            \pi_k
            \mathcal{N}(\vec{x}_n\mid\vec{\mu}_k,\Sigma_k)
        }{
            \sum_{j=1}^{K}
            \pi_j
            \mathcal{N}(\vec{x}_n\mid\vec{\mu}_j,\Sigma_j)
        }.
    \]
    \item Für GMMs aktualisiert der M-Schritt:
    \[
        \vec{\mu}_k
        =
        \frac{1}{N_k}
        \sum_{n=1}^{N}
        \gamma_{nk}\vec{x}_n,
    \]
    \[
        \Sigma_k
        =
        \frac{1}{N_k}
        \sum_{n=1}^{N}
        \gamma_{nk}
        (\vec{x}_n-\vec{\mu}_k)
        (\vec{x}_n-\vec{\mu}_k)^{\top},
    \]
    \[
        \pi_k
        =
        \frac{N_k}{N}.
    \]
    \item EM findet im Allgemeinen nur lokale Optima.
    \item Initialisierung, numerische Stabilität und Regularisierung sind praktisch wichtig.
\end{itemize}

\section{Typische Prüfungsfragen}

\begin{examtipbox}
Mögliche Prüfungsfragen zu diesem Kapitel:

\begin{enumerate}
    \item Was ist eine latente Variable?
    \item Was ist der Unterschied zwischen beobachteten und vollständigen Daten?
    \item Was ist die beobachtete Daten-Likelihood?
    \item Was ist die vollständige Daten-Log-Likelihood?
    \item Warum ist Maximum-Likelihood mit latenten Variablen oft schwierig?
    \item Was ist die \(Q\)-Funktion?
    \item Was passiert im E-Schritt?
    \item Was passiert im M-Schritt?
    \item Leite die EM-Untergrenze mit Jensen-Ungleichung her.
    \item Wie hängt die EM-Untergrenze mit der KL-Divergenz zusammen?
    \item Warum wird die Schranke im E-Schritt eng?
    \item Warum verringert EM die Log-Likelihood nicht?
    \item Warum garantiert EM kein globales Maximum?
    \item Schreibe die vollständige Daten-Likelihood eines GMMs auf.
    \item Leite die vollständige Daten-Log-Likelihood eines GMMs her.
    \item Was sind Responsibilities?
    \item Leite die Responsibilities für ein GMM mit Bayes' Regel her.
    \item Schreibe die \(Q\)-Funktion für ein GMM auf.
    \item Leite das Update für \(\vec{\mu}_k\) her.
    \item Leite das Update für \(\Sigma_k\) her.
    \item Leite das Update für \(\pi_k\) mit Lagrange-Multiplikator her.
    \item Wie hängt EM für GMMs mit \(k\)-means zusammen?
    \item Was ist Generalized EM?
    \item Was ist Hard EM?
    \item Wie kann EM bei fehlenden Daten verwendet werden?
    \item Welche numerischen Probleme können bei EM auftreten?
    \item Warum können GMMs singuläre Lösungen haben?
\end{enumerate}
\end{examtipbox}